% GENERATED FILE. Edit canonical lessons, then run npm run generate:books.
% canonical-source-hash: fnv1a64:63983c5d7321b4ce
% canonical-lessons: BN-C525-sur, BN-C525-tal, BN-C525-svorlipi, BN-C525-gomgom-kora, BN-C525-shis-deoya

\chapter{Making Music}
\label{ch:bn-making-music}
\hlchaptermodality{525}

\begin{chapteropening}
\textbf{By the end of this chapter:} \emph{I can talk about tunes, rhythm, notation and practising singing.}

Complete the last lesson of chapter 525: I can talk about tunes, rhythm, notation and practising singing.
\end{chapteropening}

\section[\texorpdfstring{\bn{সুর} (sur)}{সুর (sur)}]{\bn{সুর} (sur) — a tune; a melody}

\label{lesson:BN-C525-sur}



\begin{quote}
Before the new one: say the Bengali for to carve, then the Bengali for to dye.
\end{quote}

\subsection*{You'll want to know: \bn{সুর}}
\textbf{\bn{সুর}} — \emph{sur} — \textquotedblleft{}a tune; a melody\textquotedblright{}.

\begin{quote}
\bn{এই} \bn{গানের} \bn{সুর} \bn{সুন্দর।}
\end{quote}

\begin{quote}
\emph{Ei gāner sur sundôr.}
\end{quote}

\begin{quote}
This song has a lovely tune.
\end{quote}

\begin{cousinweb}
\textbf{\bn{সুর}} (\emph{sur}) goes back to Sanskrit \emph{svara}, \textquotedblleft{}sound, a note\textquotedblright{} (Monier-Williams), worn down in everyday speech.
\end{cousinweb}

\begin{grammarlens}[title={What you've built}]
One more word for describing what you lived and felt.
\end{grammarlens}

\subsection*{Your turn}
\begin{itemize}
\raggedright
  \item \emph{Say it:} \emph{sur}
  \item \emph{Say it:} \emph{sur}, once more
  \item \emph{Say it:} \emph{Ei gāner sur sundôr.} — This song has a lovely tune.
  \item \emph{Recall:} say the Bengali for to carve, then the Bengali for to dye, then say \emph{sur} again
\end{itemize}

\begin{tcolorbox}[breakable,colback=teal!4,colframe=teal!35!black,title={Before you move on}]
What does \bn{সুর} mean? (\textquotedblleft{}A tune; a melody\textquotedblright{}.) Say it once more.
\end{tcolorbox}
\section[\texorpdfstring{\bn{তাল} (tāl)}{তাল (tāl)}]{\bn{তাল} (tāl) — rhythm; a beat}

\label{lesson:BN-C525-tal}



\begin{quote}
Before the new one: say the Bengali for to dye, then the Bengali for a tune.
\end{quote}

\subsection*{You'll want to know: \bn{তাল}}
\textbf{\bn{তাল}} — \emph{tāl} — \textquotedblleft{}rhythm; a beat\textquotedblright{}.

\begin{quote}
\bn{তবলা} \bn{গানের} \bn{তাল} \bn{রাখে।}
\end{quote}

\begin{quote}
\emph{Tôblā gāner tāl rākhe.}
\end{quote}

\begin{quote}
The tabla keeps the rhythm of the song.
\end{quote}

\begin{cousinweb}
\textbf{\bn{তাল}} (\emph{tāl}) is a tadbhava: CDIAL lists Bengali \emph{tāl}, \textquotedblleft{}clapping\textquotedblright{}, under Sanskrit \emph{tāḍa}, where \emph{tāla} means a clap and musical time (CDIAL 5748).
\end{cousinweb}

\begin{grammarlens}[title={What you've built}]
One more word for describing what you lived and felt.
\end{grammarlens}

\subsection*{Your turn}
\begin{itemize}
\raggedright
  \item \emph{Say it:} \emph{tāl}
  \item \emph{Say it:} \emph{tāl}, once more
  \item \emph{Say it:} \emph{Tôblā gāner tāl rākhe.} — The tabla keeps the rhythm of the song.
  \item \emph{Recall:} say the Bengali for to dye, then the Bengali for a tune, then say \emph{tāl} again
\end{itemize}

\begin{tcolorbox}[breakable,colback=teal!4,colframe=teal!35!black,title={Before you move on}]
What does \bn{তাল} mean? (\textquotedblleft{}Rhythm; a beat\textquotedblright{}.) Say it once more.
\end{tcolorbox}
\section[\texorpdfstring{\bn{স্বরলিপি} (svôrlipi)}{স্বরলিপি (svôrlipi)}]{\bn{স্বরলিপি} (svôrlipi) — musical notation}

\label{lesson:BN-C525-svorlipi}



\begin{quote}
Before the new one: say the Bengali for a tune, then the Bengali for rhythm.
\end{quote}

\subsection*{You'll want to know: \bn{স্বরলিপি}}
\textbf{\bn{স্বরলিপি}} — \emph{svôrlipi} — \textquotedblleft{}musical notation\textquotedblright{}.

\begin{quote}
\bn{মেয়েটা} \bn{খাতায়} \bn{গানের} \bn{স্বরলিপি} \bn{লিখল।}
\end{quote}

\begin{quote}
\emph{Meyeṭā khātāy gāner svôrlipi likhlô.}
\end{quote}

\begin{quote}
The girl wrote down the song's notation in a notebook.
\end{quote}

\begin{cousinweb}
\textbf{\bn{স্বরলিপি}} (\emph{svôrlipi}) joins Sanskrit \emph{svara}, \textquotedblleft{}a note\textquotedblright{}, and \emph{lipi}, \textquotedblleft{}writing, a script\textquotedblright{} (Monier-Williams): music written down.
\end{cousinweb}

\begin{grammarlens}[title={What you've built}]
One more word for describing what you lived and felt.
\end{grammarlens}

\subsection*{Your turn}
\begin{itemize}
\raggedright
  \item \emph{Say it:} \emph{svôrlipi}
  \item \emph{Say it:} \emph{svôrlipi}, once more
  \item \emph{Say it:} \emph{Meyeṭā khātāy gāner svôrlipi likhlô.} — The girl wrote down the song's notation in a notebook.
  \item \emph{Recall:} say the Bengali for a tune, then the Bengali for rhythm, then say \emph{svôrlipi} again
\end{itemize}

\begin{tcolorbox}[breakable,colback=teal!4,colframe=teal!35!black,title={Before you move on}]
What does \bn{স্বরলিপি} mean? (\textquotedblleft{}Musical notation\textquotedblright{}.) Say it once more.
\end{tcolorbox}
\section[\texorpdfstring{\bn{গমগম} \bn{করা} (gômgôm kôrā)}{গমগম করা (gômgôm kôrā)}]{\bn{গমগম} \bn{করা} (gômgôm kôrā) — to resound; to boom}

\label{lesson:BN-C525-gomgom-kora}



\begin{quote}
Before the new one: say the Bengali for rhythm, then the Bengali for musical notation.
\end{quote}

\subsection*{You'll want to know: \bn{গমগম} \bn{করা}}
\textbf{\bn{গমগম} \bn{করা}} — \emph{gômgôm kôrā} — \textquotedblleft{}to resound; to boom\textquotedblright{}.

\begin{quote}
\bn{গানের} \bn{সময়} \bn{পুরো} \bn{ঘরটা} \bn{গমগম} \bn{করছিল।}
\end{quote}

\begin{quote}
\emph{Gāner sômôy puro ghôrṭā gômgôm kôrchhilô.}
\end{quote}

\begin{quote}
During the song the whole room was resounding.
\end{quote}

Past (he/she …): \bn{সেই} \bn{ঘর} \bn{গানে} \bn{গমগম} \bn{করল}, \emph{sei ghôr gāne gômgôm kôrlô}.

\begin{cousinweb}
\textbf{\bn{গমগম} \bn{করা}} (\emph{gômgôm kôrā}) joins \bn{গমগম}, an imitative word for a deep, filling sound, to \bn{করা}, \textquotedblleft{}to do\textquotedblright{}. No dictionary we use records an older form, so its history is not settled.
\end{cousinweb}

\begin{grammarlens}[title={What you've built}]
One more word for describing what you lived and felt.
\end{grammarlens}

\subsection*{Your turn}
\begin{itemize}
\raggedright
  \item \emph{Say it:} \emph{gômgôm kôrā}
  \item \emph{Say it:} \emph{gômgôm kôrā}, once more
  \item \emph{Say it:} \emph{Gāner sômôy puro ghôrṭā gômgôm kôrchhilô.} — During the song the whole room was resounding.
  \item \emph{Recall:} say the Bengali for rhythm, then the Bengali for musical notation, then say \emph{gômgôm kôrā} again
\end{itemize}

\begin{tcolorbox}[breakable,colback=teal!4,colframe=teal!35!black,title={Before you move on}]
What does \bn{গমগম} \bn{করা} mean? (\textquotedblleft{}To resound; to boom\textquotedblright{}.) Say it once more.
\end{tcolorbox}
\section[\texorpdfstring{\bn{শিস} \bn{দেওয়া} (shis deoyā)}{শিস দেওয়া (shis deoyā)}]{\bn{শিস} \bn{দেওয়া} (shis deoyā) — to whistle}

\label{lesson:BN-C525-shis-deoya}



\begin{quote}
Before the new one: say the Bengali for musical notation, then the Bengali for to resound.
\end{quote}

\subsection*{You'll want to know: \bn{শিস} \bn{দেওয়া}}
\textbf{\bn{শিস} \bn{দেওয়া}} — \emph{shis deoyā} — \textquotedblleft{}to whistle\textquotedblright{}.

\begin{quote}
\bn{ছেলেটা} \bn{শিস} \bn{দিয়ে} \bn{একটা} \bn{গান} \bn{গাইল।}
\end{quote}

\begin{quote}
\emph{Chheleṭā shis diye ekṭā gān gāilô.}
\end{quote}

\begin{quote}
The boy whistled a song.
\end{quote}

Past (he/she …): \bn{তিনি} \bn{শিস} \bn{দিলেন}, \emph{tini shis dilen}.

\begin{cousinweb}
\textbf{\bn{শিস} \bn{দেওয়া}} (\emph{shis deoyā}) joins \bn{শিস}, probably an imitative word for a whistle, to \bn{দেওয়া}, \textquotedblleft{}to give\textquotedblright{}. No dictionary we use records an older form, so its history is not settled.
\end{cousinweb}

\begin{grammarlens}[title={What you've built}]
One more word for describing what you lived and felt.
\end{grammarlens}

\subsection*{Your turn}
\begin{itemize}
\raggedright
  \item \emph{Say it:} \emph{shis deoyā}
  \item \emph{Say it:} \emph{shis deoyā}, once more
  \item \emph{Say it:} \emph{Chheleṭā shis diye ekṭā gān gāilô.} — The boy whistled a song.
  \item \emph{Recall:} say the Bengali for musical notation, then the Bengali for to resound, then say \emph{shis deoyā} again
\end{itemize}

\begin{tcolorbox}[breakable,colback=teal!4,colframe=teal!35!black,title={Before you move on}]
What does \bn{শিস} \bn{দেওয়া} mean? (\textquotedblleft{}To whistle\textquotedblright{}.) Say it once more.
\end{tcolorbox}
